\chapter{Checking the certificates and understanding the scope}\label{ch:reproduction}

The volume theorems combine written analytic and geometric arguments with finite checks of different kinds. The long-twist proof reduces two local inequalities to exact rational certificates supplied with the paper. The quantitative circle-tree proof establishes its uniform bounds analytically; separate finite checks verify its edge formulas, face identifications and peripheral identities. This chapter distinguishes those roles and records the scope of each theorem.

\section{Three different kinds of evidence}
The manuscript contains analytic proofs, exact finite certificates, and optional finite audits. They have different logical roles. The matrix identities, support and strictness arguments, geometric incidence, complex estimates and filling analysis are proved in the text. The rational certificates are finite parts of the proof: the continuum inequalities reduce to their exact feasibility, coverage and supporting-plane checks. Finite matrix or numerical asymptotic audits provide independent checks of formulas and normalizations, but do not establish a theorem by extrapolation.

The proof is given in the chapters and appendices of this manuscript. The repository also supplies the exact certificate data and verification programs. Those programs check the finite inequalities described below; they do not formally verify the analytic or geometric arguments.

\section{Replaying the exact certificates}
From the repository root, run
\begin{verbatim}
python3 repro/verify.py --output repro/results
\end{verbatim}
The serial runner uses the standard Python library. It replays the two-leaf block certificate, the attachment square certificate and the optional endpoint certificate, compares their exact results with the recorded data, and checks the elementary rational leaf and isolated-three-leaf margins. Its report records the source and data hashes. Candidate generation and numerical optimization are not run by verification.

The principal proof files are
\begin{center}\small
\begin{tabular}{p{0.44\textwidth}p{0.46\textwidth}}\toprule
Certificate & Mathematical coverage\\\midrule
\path{verify_conway_two_leaf_block_alias_certificate.py}
& Sixteen slope intervals; both alias types; 288 corner/vertex inequalities.\\[5pt]
\path{verify_conway_single_leaf_attachment_certificate.py}
& Exact 55-rectangle cover; 880 corner/vertex inequalities on the full child/parent slope square.\\[5pt]
Endpoint certificate in the runner's manifest
& Optional stronger margin at the monotone-response endpoint; not needed for the main theorem.\\\bottomrule
\end{tabular}
\end{center}
All paths are resolved relative to the repository or script directory, rather than an author's temporary working directory. The runner's proof-input guide lists the written dependencies separately from executable checks.

\section{Checking the quantitative tree formulas}
The elementary edge equations and complex-strip estimates used by Chapter~\ref{ch:quantitative-circle-trees} have a portable reproduction command:
\begin{verbatim}
python3 repro/run_research_checks.py curvature-fillings
\end{verbatim}
The runner checks exact symbolic edge identities, finite samples of the stated complex inequality, labeled octahedral attachments and the full signed peripheral matrix words. Its output separates those computations from the uniform analytic and geometric proofs. Appendix~\ref{app:quantitative-sums} gives the full quantitative finite-sum proof, including both uniform Newton estimates, the weighted contraction and the charge-tail bounds. Appendix~\ref{app:quantitative-geometry} gives the complete ordered attachment, positive angular charts, primitive-period completion, volume calibration and preferred-cusp argument. The linked versions are optional separately available copies of those proofs. Such finite checks can detect sign or indexing errors; they do not replace the written proofs.

\section{The planar terminal decoration}
The finite identities for Chapter~\ref{ch:balanced-terminal-fans} have one command:
\begin{verbatim}
python3 repro/run_research_checks.py balanced-terminal-fans
\end{verbatim}
The script uses exact integer and rational arithmetic. It checks the
closed fan recurrence, the rational balance margin, the oriented
triangle subdivision and its new ordered quadrilateral, the matching
and even-level face parity, and the endpoint pairing of even stacks of
odd terminal tangles. It does not infer an asymptotic theorem from
sampled quantum values. Appendix~\ref{app:balanced-terminal-fans}
proves the geometric realization, full finite source, uniform
Fourier estimates and primitive completion.

\section{The directed criterion and unmarked junctions}

The exact rational comparisons for
Chapter~\ref{ch:directed-circle-messages} have one command:
\begin{verbatim}
python3 repro/run_research_checks.py directed-tree-fillings
\end{verbatim}
This program uses only Python's standard library. It verifies the
constant inequalities for the uniform twist bound and the star with
central power two, using rational enclosures of all quantities.
The uniform conditional kernel estimate, induction over arbitrary
fixed trees, noncancellation and geometric interpretation are written
proofs in Appendix~\ref{app:directed-circle-messages}.

For the local algebra of the unmarked-junction extension, run
\begin{verbatim}
python3 repro/run_research_checks.py strict-tree-geometry
\end{verbatim}
This program requires SymPy. It checks the white-map normalizations,
primitive disk-longitude trace, incident radius formulas and the
half-angle fold identity. Appendix~\ref{app:strict-tree-geometric-characteristic}
proves the based boundary-map comparison, stationary elimination and
full representation assembly. Those global statements are not
deduced from the finite matrix checks.

\section{Building the manuscript and illustrated guide}
The TeX source is in \sourcefile{manuscript/}; its native illustrations are in \sourcefile{manuscript/figures/}. A TeX Live or MacTeX installation with \sourcefile{latexmk} builds the paper:
\begin{verbatim}
python3 tools/build_manuscript.py
\end{verbatim}
The output is \sourcefile{manuscript/build/main.pdf}. The build manifest records the hashes of all TeX and bibliography inputs and the generated PDF. It is a typesetting record, not a mathematical certificate.

The illustrated guide is built separately:
\begin{verbatim}
python3 tools/build_book.py
python3 tools/check_book.py
\end{verbatim}
It is written to \sourcefile{docs/} and links the same authored PDF. Documentation builds do not execute research scripts. The guide also links the broader research companions and states their hypotheses separately from the two theorem classes developed here.

\section{What has and has not been proved}
Theorem~\ref{thm:main} is uniform in the level for each fixed finite completed expression and each fixed sufficiently long signed filling tuple. It includes both parity sequences and composite levels. Its threshold can depend on the expression, its complete geometry and its retained half-twist pattern, but imposes no bound on ratios between the long slopes.

Theorem~\ref{thm:circle-quantitative} uses a different sufficient hypothesis: the literal ordered circle tree and the explicit sum $\sum_vd_v/(d_v^2+n_v^2)\le2^{-25}$. The full diagram, all degrees and all integer powers are fixed before the limit. The theorem gives actual hyperbolicity and volume equality, a uniform microscopic meridian profile with its exact scalar, and the preferred-cusp correction. It permits short coefficients under this condition; it does not prove a statement about arbitrary short or rational fillings.

Both results keep the original strand order, knot-pairing test, disk framings and unfilled main cusp. The octahedral attachment is proved by actual labeled cuts. A scalar equality in $\cH_N$ is not used as a tangle isotopy. The continued boundary word used for the preferred-cusp calculation becomes the preferred longitude of the smooth filled knot only at the endpoint of the geometric path.

These are sufficient construction hypotheses, not a classification of all arborescent knots. The completed-box condition remains part of Theorem~\ref{thm:main}; the quarter-index states in circle trees are treated by the separate proof in Chapter~\ref{ch:quantitative-circle-trees}. More general strict ordered Conway trees have a long-twist theorem in the linked research companions, with their own full source and marked-geometry proofs. No general short-filling theorem follows merely from the existence of complete augmented cells or from an upper bound on the absolute values of a larger sum.

\section{The geometric longitude calculation}
The symbolic identities in the geometric characteristic chapter have one
reproduction command:
\begin{verbatim}
python3 repro/run_research_checks.py geometric-characteristic
\end{verbatim}
This group requires SymPy. It verifies the nonparabolic face maps,
character matrices, half-twist line transports, and the differential
and discriminant formulas. The topological gluing, boundary pairing,
algebraic constants and regularity arguments remain the written proofs
of Appendices~\ref{app:filtered-characteristic} and
\ref{app:geometric-characteristic}.


\section{Uniform bounds for inserting twists}
Appendix~\ref{app:uniform-filling-bound} proves a level-independent
bound for each original filling functional, measured against the
supremum of the shared spectral polynomial. Its finite identities
can be checked with one command:
\begin{verbatim}
python3 repro/run_research_checks.py filling-supremum
\end{verbatim}
This group requires SymPy and NumPy. It checks the derivative
interpolation and Chebyshev residues exactly, then compares the
functional with independently constructed signed crossing matrices
over finite fields. The shared exponential estimate remains a
separate hypothesis; these checks add no new volume class.
